\documentclass[12pt,a4paper]{amsart}

\usepackage{amsmath,amssymb,amsthm}
\usepackage{mathtools}
\usepackage{hyperref}
\usepackage{geometry}
\usepackage[ngerman]{babel}
\geometry{margin=2.5cm}

% -------------------------------------------------------
\newtheorem{theorem}{Satz}[section]
\newtheorem{lemma}[theorem]{Lemma}
\newtheorem{corollary}[theorem]{Korollar}
\newtheorem{proposition}[theorem]{Proposition}
\theoremstyle{definition}
\newtheorem{definition}[theorem]{Definition}
\newtheorem{remark}[theorem]{Bemerkung}

\newcommand{\Z}{\mathbb{Z}}
\newcommand{\euler}{\varphi}

% -------------------------------------------------------
\begin{document}
% -------------------------------------------------------

\title{Kein Semiprim erfüllt Lehmers Totient-Bedingung:\\ Ein arithmetischer Beweis}

\author{Kurt Ingwer}
\address{Specialist Mathematics Project, Version Build 44}
\email{specialist-maths@localhost}
\date{\today}

\begin{abstract}
Lehmers Totient-Problem fragt, ob $\euler(n) \mid (n-1)$ die Primalität von~$n$ erzwingt.
Wir beweisen, dass kein Semiprim $n = p \cdot q$ mit $p < q$ prim diese Bedingung erfüllt.
Die Schlüsselidentität ist $pq - 1 \equiv p + q - 2 \pmod{(p-1)(q-1)}$: Da
$(p-1)(q-1) > p + q - 2$ für $p \ge 2$, $q > p$, kann der Quotient $(p-1)(q-1)$
die Zahl $pq - 1$ nicht teilen. Dies schließt Semiprimen aus Lehmers Problem aus
und zeigt, dass jede Lehmer-Zahl (falls eine existiert) mindestens drei verschiedene
Primfaktoren haben muss.
\end{abstract}

\maketitle

% -------------------------------------------------------
\section{Einleitung}
% -------------------------------------------------------

Lehmer \cite{Lehmer1932} fragte, ob $\euler(n) \mid (n-1)$ die Primalität von~$n$ erzwingt.
Nachdem gezeigt wurde, dass jede Lösung quadratfrei sein muss \cite{Ingwer2026qfrei},
ist der nächste Schritt, Produkte genau zweier Primzahlen auszuschließen.

Für ein quadratfreies $n = p_1 \cdots p_k$ lautet die Lehmer-Bedingung
\[
  (p_1-1)(p_2-1)\cdots(p_k-1) \mid p_1 p_2 \cdots p_k - 1.
\]
Für $k = 2$ mit $p < q$ prim ist dies $(p-1)(q-1) \mid pq - 1$.

% -------------------------------------------------------
\section{Hauptresultat}
% -------------------------------------------------------

\begin{theorem}
\label{thm:keinsemi}
Es existiert kein Semiprim $n = p \cdot q$ mit $p < q$ prim, das
$\euler(n) \mid (n-1)$ erfüllt.
\end{theorem}

\begin{proof}
Die Bedingung $\euler(pq) \mid (pq - 1)$ lautet $(p-1)(q-1) \mid (pq - 1)$.

Wir verwenden die Schlüsselidentität:
\begin{equation}
  pq - 1 = (p-1)(q-1) + (p-1) + (q-1).
  \label{eq:ident}
\end{equation}

\begin{proof}[Verifikation von \eqref{eq:ident}]
$(p-1)(q-1) + (p-1) + (q-1) = pq - p - q + 1 + p - 1 + q - 1 = pq - 1$. \checkmark
\end{proof}

Aus~\eqref{eq:ident} folgt:
\[
  pq - 1 \equiv (p-1) + (q-1) = p + q - 2 \pmod{(p-1)(q-1)}.
\]
Damit $(p-1)(q-1) \mid (pq-1)$ gilt, muss $(p-1)(q-1) \mid (p + q - 2)$ gelten.

Wir zeigen, dass dies unmöglich ist:

\textbf{Fall $p = 2$:} Die Bedingung lautet $(q-1) \mid q$.
Da $q = (q-1)+1$, gilt $q \equiv 1 \pmod{q-1}$, also $(q-1) \mid 1$,
d.\,h.\ $q-1 = 1$, also $q = 2 = p$.
Dies widerspricht $q > p$.

\textbf{Fall $p \ge 3$:} Da $p$ und $q$ ungerade Primzahlen mit $q > p$ sind,
gilt $q \ge p + 2$, also $q - 2 \ge p \ge 3$ und $p - 2 \ge 1$. Damit:
\[
  (p-1)(q-1) - (p+q-2) = (p-2)(q-2) - 1 \ge 1 \cdot p - 1 = p - 1 \ge 2 > 0.
\]
Also $(p-1)(q-1) > p + q - 2 \ge 1$, und daher $(p-1)(q-1) \nmid (p+q-2)$.

In beiden Fällen gilt $(p-1)(q-1) \nmid (pq-1)$ — Widerspruch.
\end{proof}

% -------------------------------------------------------
\section{Die Schlüsselidentität im Kontext}
% -------------------------------------------------------

\begin{remark}
Identität~\eqref{eq:ident} verallgemeinert sich: Für $n = p_1 \cdots p_k$ quadratfrei gilt
\[
  n - 1 \equiv \sum_{i=1}^{k}(p_i - 1) \pmod{\euler(n)}.
\]
Dies ist die modulare Reduktion, die der obige Beweis für $k = 2$ nutzt.
Für größeres~$k$ wird die Analyse aufwändiger, aber das Grundprinzip bleibt dasselbe.
\end{remark}

\begin{remark}
Die Ungleichung $(p-1)(q-1) > p + q - 2$ ist äquivalent zu $(p-2)(q-2) > 1$,
was für alle Primpaare $p \ge 3$, $q > p$ gilt (da $p - 2 \ge 1$ und $q - 2 \ge p - 1 \ge 2$).
\end{remark}

\begin{corollary}
Jede Lehmer-Zahl hat mindestens drei verschiedene Primfaktoren.
\end{corollary}

\begin{proof}
Quadratfreiheit \cite{Ingwer2026qfrei} schließt Primzahlpotenzen aus.
Satz~\ref{thm:keinsemi} schließt Produkte zweier verschiedener Primzahlen aus.
Also werden mindestens drei benötigt.
\end{proof}

% -------------------------------------------------------
\begin{thebibliography}{10}

\bibitem{Lehmer1932}
D.~H.~Lehmer,
\emph{On Euler's totient function},
Bull.\ Amer.\ Math.\ Soc.\ \textbf{38} (1932), 745--751.

\bibitem{Cohen1980}
G.~L.~Cohen und P.~Hagis,
\emph{On the number of prime factors of $n$ if $\phi(n) \mid (n-1)$},
Nieuw Arch.\ Wisk.\ (3) \textbf{28} (1980), 177--185.

\bibitem{Ingwer2026qfrei}
K.~Ingwer,
\emph{Jede Lösung von Lehmers Totient-Problem ist quadratfrei},
Specialist Mathematics Project (2026).

\bibitem{Ingwer2026lehmer3}
K.~Ingwer,
\emph{Lehmers Totient-Problem für Produkte dreier Primzahlen},
Specialist Mathematics Project (2026).

\end{thebibliography}

\end{document}
